% !PRINTED 449
\weberpage{449}
%<<CONT
so ist wegen 5. $N(p) = p^{2h}$, und folglich
\[
2h = r\varepsilon.
\]

Die Anzahl der Primfaktoren, in die $p$ im Körper $\mathfrak{K}$ zerfällt,
ist also gleich $2h/\varepsilon$.

Ist ferner $(D, p) = -1$, so sind die Primfaktoren von $p$ vom zweiten Grade, und ihre Anzahl ist also $= h$.

9. Eine interessante Folgerung ziehen wir noch aus diesen Resultaten. Es sei $(D, p) = +1$ und $\varepsilon$ der Index von $p$. Dann ist $4p^{\varepsilon}$ durch die Form $x^2 - Dy^2$ darstellbar und
\weq{23}{p^{\varepsilon} = \mu\mu',}
wenn zur Abkürzung
\weq{24}{\mu = \frac{x + y\sqrt{D}}{2}, \quad \mu' = \frac{x - y\sqrt{D}}{2}}
gesetzt wird. Wir wissen nun, daß $p$ in lauter verschiedene Primideale vom Grade $\varepsilon$ zerfällt. Da $p$ ungerade ist, und $x, y$ höchstens den gemeinschaftlichen Teiler 2 haben, so haben $\mu, \mu'$ keinen gemeinsamen Primfaktor in $\mathfrak{K}$. Jedem Primideal $\mathfrak{P}$, das in $\mu$ aufgeht, entspricht ein davon verschiedenes (konjugiertes) Primideal $\mathfrak{P}'$, das aus $\mathfrak{P}$ dadurch entsteht, daß man für alle Zahlen von $\mathfrak{P}$ die konjugiert imaginären Zahlen setzt. Es ist also $\mu$ nicht nur durch $\mathfrak{P}$, sondern durch $\mathfrak{P}^{\varepsilon}$ teilbar.

Nach \S\,121 ist
\weq{25}{\left(\frac{\eta(\omega_1)}{\eta(\omega)}\right)^2 \left(\frac{\eta(\omega_2)}{\eta(\omega_1)}\right)^2 \cdots \left(\frac{\eta(\omega_{\varepsilon})}{\eta(\omega_{\varepsilon-1})}\right)^2 = \varrho\mu,}
und nach \S\,72, 4 genügen die Faktoren $P^2$ auf der linken Seite dieses Ausdrucks einer Gleichung $P^2\varphi(P^2, \gamma_2, \gamma_3) = p$, worin $\varphi(P^2, \gamma_2, \gamma_3)$ eine ganze algebraische Zahl ist.

Daraus folgt aber, daß $P^2$ nicht durch eine höhere als die erste Potenz von $\mathfrak{P}$ teilbar sein kann, während doch $\mu$ durch $\mathfrak{P}^{\varepsilon}$ teilbar ist. Wir schließen also aus (25), daß \gesperrt{jede} der Zahlen $P^2$ durch die erste Potenz von $\mathfrak{P}$ teilbar sein muß, daß mithin alle diese Zahlen \gesperrt{assoziiert sind}. Bezeichnet also $\varrho$ irgend eine algebraische Einheit, so ist
\[
\frac{x + y\sqrt{-m}}{2} = \varrho P^{2\varepsilon},
\]
also ist $\mu$ wirklich als $\varepsilon$te Potenz einer im Körper $\mathfrak{K}$ existierenden Zahl dargestellt, und $P^2$ selbst ist eine Darstellung eines Primideals im Körper $\Omega$ durch ein \gesperrt{Hauptideal} im Körper $\mathfrak{K}$.
